\documentclass[10pt, a4paper]{article}
\usepackage{preamble}

\title{Probability II \\
    \large Cheat Sheet}
\author{Luke Phillips}
\date{March 2026}

\begin{document}

\maketitle

\newpage

\section{Probability II -
    Probability Spaces}

\begin{proposition}[Boole's inequality]
    For events $A_1, A_2, \dotsc$,
    \[
    \P\left(\bigcup_{k = 1}^{\infty}A_k\right) \leq \sum_{k = 1}^{\infty}\P(A_k).
    \]
\end{proposition}

\begin{proposition}[Markov's inequality]
    $X \geq 0$ non-negative random variable then for every $a > 0$,
    \[
    \P(X \geq a) \leq \frac{\E(X)}{a}.
    \]
\end{proposition}

\begin{proposition}[Chebyshev's inequality]
    $X$ random variable then for every $a > 0$,
    \[
    \P(|X - \E(X)| \geq a) \leq \frac{\Var(X)}{a ^ 2}.
    \]
\end{proposition}

\begin{definition}[Sigma-field]
    $\mathcal{F}$ is a sigma field if
    \begin{enumerate}[label = (\roman*)]
        \item $\emptyset \in \mathcal{F}$.

        \item If $A \in \mathcal{F}$,
        then $A ^ c \in \mathcal{F}$.

        \item If $A_1, A_2, \dotsc \in \mathcal{F}$,
        then $\bigcup_{k = 1}^{\infty}A_k \in \mathcal{F}$.
    \end{enumerate}
\end{definition}

\begin{lemma}
    If $(A_n)_{n \geq 1}$ is increasing with $A \coloneqq \lim_{n}A_n \equiv \bigcup_{n \geq 1}A_n$,
    then
    \[
    \P(A) = \P\left(\lim_{n \to \infty}A_n\right) = \lim_{n \to \infty}\P(A_n).
    \]
    If $(A_n)_{n \geq 1}$ is a decreasing sequence with $A \coloneqq \lim_{n}A_n \equiv \bigcap_{n \geq 1}A_n$,
    then
    \[
    \P(A) = \P\left(\lim_{n \to \infty}A_n\right) = \lim_{n \to \infty}\P(A_n).
    \]
\end{lemma}


\end{document}